\documentclass[11pt,a4paper]{article}

% ===== XeLaTeX + Korean =====
\usepackage{kotex}
\usepackage{fontspec}
\setmainfont{Times New Roman}
\setmainhangulfont{AppleMyungjo}
\setsanshangulfont{Apple SD Gothic Neo}

% ===== Layout =====
\usepackage[a4paper,top=18mm,bottom=18mm,left=20mm,right=20mm]{geometry}
\usepackage{fancyhdr}
\setlength{\headheight}{24pt}
\usepackage{enumitem}
\usepackage{mdframed}
\usepackage{array}

\setlength{\parindent}{1em}
\linespread{1.18}

% ===== Header / Footer =====
\pagestyle{fancy}
\fancyhf{}
\renewcommand{\headrulewidth}{0.6pt}
\fancyhead[L]{\sffamily\bfseries 2026 고1 영어 변형 — 지문 31}
\fancyhead[R]{\sffamily 변화와 Heraclitus}
\renewcommand{\footrulewidth}{0pt}
\fancyfoot[C]{\framebox[16mm][c]{\thepage}}

% ===== helpers =====
\newcommand{\qno}[1]{\par\vspace{8pt}\noindent{\sffamily\bfseries #1.}\ }
\newcommand{\instr}[1]{{\sffamily #1}\par\vspace{3pt}}
\newlist{choices}{enumerate}{1}
\setlist[choices]{label=,leftmargin=1.5em,itemsep=2pt,topsep=3pt,parsep=0pt}

\newmdenv[linewidth=0.6pt,roundcorner=3pt,innertopmargin=7pt,innerbottommargin=7pt,
          innerleftmargin=9pt,innerrightmargin=9pt,skipabove=5pt,skipbelow=5pt]{pbox}
\newcommand{\nemo}[1]{\fbox{\,#1\,}}

\begin{document}

\begin{center}
{\fontsize{15}{17}\selectfont\sffamily\bfseries [지문 31] 다음 글을 읽고 물음에 답하시오.}
\end{center}
\vspace{4pt}

% ===== 공통 지문 (빈칸: imperceptibly 자리, passive observers 복원) =====
\begin{pbox}
Our world is not just interconnected, but ever-changing, even if we can't sense it.
While you're reading this, you're changing.
You're aging (a minute amount, thankfully), but the neural$^{*}$ networks in your brain
are also \rule{3.5cm}{0.4pt} changing as you perceive each word.
Crucially, even when we're seemingly not doing anything of note, events are taking
place outside your immediate surroundings that will change your life in the future,
though you won't realize it yet.
Heraclitus, the ancient Greek philosopher, rightly pointed out, ``No man ever steps
in the same river twice. For it's not the same river and he's not the same man.''
To that, Cratylus, a student of Heraclitus, added that we are not mere passive observers.
When you step in a river, you change it.
Nothing is static$^{***}$.
Even microscopic changes add up over time.\par\vspace{4pt}
\noindent\footnotesize
$^{*}$ neural: 신경(계)의\quad
$^{**}$ imperceptibly: 감지할 수 없게\quad
$^{***}$ static: 정적(靜的)인
\end{pbox}

% ===================== Q1 빈칸추론 =====================
\qno{1}\instr{다음 빈칸에 들어갈 말로 가장 적절한 것은?}
\begin{choices}
\item ① rapidly
\item ② intentionally
\item ③ subtly
\item ④ visibly
\item ⑤ permanently
\end{choices}

% ===================== Q2 영어 주제 =====================
\qno{2}\instr{윗글의 주제로 가장 적절한 것은?}
\begin{choices}
\item ① the philosophical debate between Heraclitus and Cratylus on identity
\item ② how neural networks process language and enable perception
\item ③ the importance of mindfulness in noticing life's changes
\item ④ the continuous and often unnoticed nature of change in our world
\item ⑤ ways in which modern science confirms ancient Greek philosophy
\end{choices}

% ===================== Q3 서술형 해석 =====================
\qno{3}\instr{[서술형] 다음 문장을 우리말로 해석하시오.}
\vspace{2pt}\par\noindent
\hspace{1em}\textit{Crucially, even when we're seemingly not doing anything of note, events are taking place outside your immediate surroundings that will change your life in the future, though you won't realize it yet.}
\par\vspace{6pt}\noindent
$\Rightarrow$ \rule{\linewidth}{0.4pt}
\par\vspace{4pt}\noindent
\rule{\linewidth}{0.4pt}

% ===================== Q4 =====================
\qno{4}\instr{Heraclitus의 강물 비유가 윗글에서 의미하는 바로 가장 적절한 것은?}
\begin{choices}
\item ① Human identity remains fixed even when surroundings change.
\item ② Only large historical events are capable of changing our lives.
\item ③ Both the world and the person experiencing it are continuously changing.
\item ④ Natural environments change, but human minds are mostly static.
\item ⑤ Change matters only when we consciously notice it happening.
\end{choices}

\end{document}
